\section{Both conductor-sixteen Tate residues in the actual arithmetic window}
\label{sec:conductor16-two-residue-window}

\paragraph{Status and dependencies.}
This is a conditional-interface proof using the original strict statements
AC-G142 and AC-G144 and the already established comparison in AC-G143.
Their class certificates are AC-G142-class and AC-G144-class.
The exact dependency-review status is recorded in the accompanying events.
The remaining arithmetic bound and mixed theta realization are not proved.
No new coefficient-prime uniformity is inferred from these statements at $269$.

Put $p=269$, $K=\mathbb Q(\omega)$, and $t^3=p$ with $t$ real.
Let $\chi$ be the fixed cubic character of $K(t)/K$, and let $\phi$
be the arithmetic quartic character modulo sixteen with
\[
 \ker\phi=\{1,7\},\qquad \phi(13)=i,\qquad
 \phi(5)=-i,\qquad \phi(-1)=-1.
 \tag{CWU.1}
\]
Changing $i$ interchanges $\phi$ and $\phi^{-1}$ and does not change
the character set used below.  The notation $c_p$ denotes the actual
cyclotomic character, with the same convention as in AC-K25.

For every source page, every $0<\varepsilon<1/256$, and its prescribed
pair $q,d$, retain the actual fields
\[
 \begin{split}
 F_u&=\mathbb Q(\zeta_{2^{h(u)+2}},\zeta_7,
                    \sqrt{-q},\sqrt{-d},\omega),\\
 H_u&=F_u(t),\qquad G_u=\operatorname{Gal}(H_u/K),\\
 u^{2+8\varepsilon}&\leq2^{h(u)}\leq2u^{2+8\varepsilon}.
 \end{split}
 \tag{CWU.2}
\]
Write $D_u=|G_u|$ and $T=T(u)=\lfloor u^{399}\rfloor$.
Use exactly the rings $O_\ell$, fraction fields $E_\ell$ and degrees
$b_\ell$ in (FC.1), and exactly the index sum $I_{399}$ in (TC.1).
For $u\geq2$ define
\[
 \begin{gathered}
 \mathscr Z_u^{\cup}=
 \{(p,\chi^{-j}\phi^{-\sigma},e):
 j\in\{1,2\},\ \sigma\in\{1,-1\},\ 1\leq e<T,\\
 \hfill e\equiv133\text{ or }267\pmod{268}\},\\
 I_{399}^{\circ\circ}(H_u,u)=
 \sum_{(\ell,\eta,e)\notin\mathscr Z_u^{\cup}}
       \frac{\log\#k_{O_\ell}}{b_\ell}
                              \max(0,i_{\ell,\eta,e}),
 \end{gathered}
 \tag{CWU.3}
\]
where the last sum has all the original ranges of (TC.1).
Only the indicated labels are removed.  In particular the quadratic
character $\tau_p$ appearing in the residual splitting field of
AC-G144 is not inserted into $\eta$ or into the host $H_u$.

\begin{theorem}[AC-G145-window: the union of the two actual Tate residues]
\label{AC-G145-window}
For every actual host (CWU.2), uniformly in its allowed page and pair,
all labels in $\mathscr Z_u^{\cup}$ are actual labels of AC-K25 and
their strict groups (FC.2) vanish on their original coefficient rings.
The number of omitted labels is exactly
\[
        |\mathscr Z_u^{\cup}|=4\left\lfloor\frac{T(u)}{134}\right\rfloor.
 \tag{CWU.4}
\]
Their contribution to (FC.15) is zero.  Their positive index costs
are retained through the single bound
\[
 \begin{split}
 0\leq I_{399}(H_u,u)-I_{399}^{\circ\circ}(H_u,u)
 \leq{}&8D_u^2T(u)^2\log(8T(u)D_u)\\
 &+\tfrac12D_u\mathcal R_uT(u)^2
                         +2D_ur_uc_0T(u)\\
 ={}&O_\varepsilon\bigl(u^{802+16\varepsilon}(\log u)^2\bigr).
 \end{split}
 \tag{CWU.5}
\]
Here $\mathcal R_u,r_u,c_0$ have exactly their meanings in (FL.1).
The majorant is not doubled when the two residues are combined.

Let $I_{399}^{\circ}$ denote the already reduced sum of AC-G143.
Then
\[
       I_{399}^{\circ\circ}\leq I_{399}^{\circ}\leq I_{399}.
 \tag{CWU.6}
\]
The common-window premise (TC.3), or its AC-G143 version, can be
replaced by the identical premise with $I_{399}^{\circ\circ}$,
retaining positive relative slack and every simultaneous quantifier.
The old premise implies this one with the same parameters.
Conversely this premise implies the full old premise after increasing
only the initial threshold and replacing the slack by half its value.
Thus it suffices for the same conditional conclusion of AC-K35 with
the same unproved $\mathcal R_{\rm cub}$ and all other retained inputs.
No upper bound for the remaining sum is asserted.
\end{theorem}

\begin{proof}
\proofstep{1. The actual finite characters and original coefficients.}
Since $u\geq2$, the lower page inequality implies $h(u)\geq2$.
Thus $\mathbb Q(\zeta_{16})\subset F_u$ and $K(t)\subset H_u$.
The characters $\phi|_{G_K}$ and $\chi$ therefore factor through
$G_u$.  The field $F_u/\mathbb Q$ is abelian, whereas $K(t)/\mathbb Q$
is the nonabelian $S_3$ splitting field.  The actual disjointness
$K(t)\cap F_u=K$ follows, so restriction to the cubic subgroup
distinguishes $j$, and the quartic quotient distinguishes $\sigma$.
There are four distinct finite characters.

The exponent $N=3\cdot2^{h(u)}$ of $G_u$ contains both orders three
and four, and is prime to $p$.  The original local ring $O_p$ in
(FC.1) therefore contains their values.  Its actual degree $b_p$
is retained; neither a degree-one assumption nor a coefficient
extension is made.  The coefficient (FC.2) is exactly
\[
 (O_p\otimes_{\mathbb Z_p}\mathbb Q_p/\mathbb Z_p)
       (\chi^{-j}\phi^{-\sigma})(-e)
       =(E_p/O_p)(\chi^{-j}\phi^{-\sigma}c_p^{-e}).
 \tag{CWU.7}
\]
Apply AC-G142 with its finite exponent zero and continuous character
one when $e\equiv267$, and AC-G144 with $m=0$, $\kappa=1$ when
$e\equiv133$.  Both statements hold on this given $O_p$ without an
absolute ramification restriction.  In the second application the
quadratic cyclotomic character is used only to identify the residual
Kummer splitting field.  Formula (CWU.7) is the same coefficient in
both applications and involves no new finite $p$-conductor label.

\proofstep{2. Exact strict local conditions and the count.}
Both strict theorems use zero in full $H^1(G_{K_v},A)$ at $v\mid p$
and zero in $H^1(I_v,A)$ at every other finite place.
These are the conditions in (FC.2).  To check the actual set
$\Sigma_p$, take a fixed minimal set containing the infinite place
and places over $2,3,p$.  At each added prime over $7,q,d$, the
coefficient (CWU.7) is unramified.  A strict cocycle restricts to a
coboundary on its inertia.  That inertia acts trivially, so the
coboundary is zero.  The cocycle identity gives zero on the closed
normal subgroup generated by these added inertia groups, which also
acts trivially on the coefficient.  The cocycle descends to the
smaller Galois group.  Conversely inflation preserves all the
conditions, and inflation is injective by inflation--restriction.
This verifies equality of the strict kernels on the actual sets.
Alternatively their inflations lie in the global strict groups
already shown to be zero by AC-G142 and AC-G144.

The two progressions form the single progression $e=134n-1$.
For a positive integer $T$, the inequalities $1\leq134n-1<T$
are equivalent to $1\leq n\leq\lfloor T/134\rfloor$.
Multiplying by the four distinct characters proves (CWU.4).
More precisely the new progression contributes
\[
 4\left\lfloor\frac{T+134}{268}\right\rfloor,
 \qquad
 \left\lfloor\frac{T}{268}\right\rfloor+
 \left\lfloor\frac{T+134}{268}\right\rfloor
                 =\left\lfloor\frac{T}{134}\right\rfloor.
 \tag{CWU.8}
\]
There is no overlap between the two progressions.  The vanishing
of these strict groups makes exactly their terms in (FC.15) zero.

\proofstep{3. Bound the union once, with the actual positive indices.}
For one removed label put
$C=R\Gamma(G_{K,\Sigma_p},O_p(\eta^{-1})(e+1))$,
using the actual codifferent-dual identification in (FL.2).
The derived determinant inequality (EU.2) gives
\[
 \frac{\log\#k_{O_p}}{b_p}\max(0,i_{p,\eta,e})
 \leq\frac{\log\#H^2(C)}{b_p}
 =\frac{\left|\log\#H^2(C)
                  -\log\#\mathcal S_{p,\eta}(e)\right|}{b_p}.
 \tag{CWU.9}
\]
The last equality uses strict vanishing and does not identify the
two cohomology groups.  For each fixed $e$, the nonnegative sum
over the omitted characters is a subsum of the absolute bound
(FL.4).  Sum that bound once over all $1\leq e<T$.
The estimates
\[
 \sum_{1\leq e<T}e\leq T^2/2,\qquad
 \log(8eD_u)\leq\log(8TD_u),\qquad T-1\leq T
\]
give the displayed three terms of (CWU.5).  This is a bound on the
union, not the sum of two separately maximized full errors.
It does not assert that any of the removed positive indices is zero.

The actual host estimates used in AC-G143 are
\[
 D_u\leq288u^{2+8\varepsilon},\qquad r_u\leq12,
 \qquad\mathcal R_u=O_\varepsilon
                     (u^{\varepsilon/5}(\log u)^2).
 \tag{CWU.10}
\]
They are uniform over the prescribed pair and page.  The first
term of (CWU.5) has exponent $802+16\varepsilon$; the second has
exponent $800+41\varepsilon/5$, and the third is smaller.
Consequently the entire error is
$o(\exp(u^\varepsilon/2))$ for each fixed positive $\varepsilon$.
Nonnegativity proves (CWU.6), and also bounds
$I_{399}^{\circ}-I_{399}^{\circ\circ}$ by that same error.

\proofstep{4. Preserve the single common window and every other input.}
Fix the same finite datum, page and $0<\varepsilon<1/256$ as in
AC-K35.  Suppose the new premise supplies $J\geq1$, $X_0$ and
$s\in(0,1)$ such that for every $X\geq X_0$ and each prescribed
$\operatorname{WinPair}(\varepsilon,X^{1/J},q,d)$ there is one
$W$, with $X^{1/J}\leq W$ and $W^2\leq X$, satisfying
\[
 I_{399}^{\circ\circ}(H_u,u)
          \leq(1-s)\exp(u^\varepsilon/2)
                  \quad\text{for every }W\leq u\leq W^2.
 \tag{CWU.11}
\]
Choose $U_0\geq2$ so that the uniform error in (CWU.5) is at most
$(s/2)\exp(u^\varepsilon/2)$ for every $u\geq U_0$.
This choice depends only on the same datum, page, $\varepsilon$
and $s$, not on $q,d$, a coefficient prime, character, weight,
or an independently selected window.  Replace $X_0$ by
$\max(X_0,U_0^J)$.  Then the very same $W$ lies above $U_0$ and
\[
 I_{399}(H_u,u)\leq(1-s/2)\exp(u^\varepsilon/2)
                       \quad(W\leq u\leq W^2).
 \tag{CWU.12}
\]
For a fixed finite list of simultaneous fields or labels, take the
maximum of their uniform error constants.  This retains that single
$W$ and never intersects separately chosen windows.
Conversely either previous premise implies (CWU.11) by (CWU.6)
with its original slack and parameters.  The little-o versions are
also equivalent.  Applying AC-K35 to (CWU.12) retains its actual
integer $R.t$, source host, family transport, all unremoved labels,
and the unproved mixed theta/Selmer and analytic realization input
$\mathcal R_{\rm cub}$.  This proves the conditional implication.
\end{proof}

\paragraph{Exact reduction of the existing arithmetic gate.}
The already reduced AC-G143 sum still contained the four actual
characters at every weight $e\equiv133\bmod268$.
For every $T\geq134$ the new omitted label set strictly contains
the old one, by (CWU.8); these are infinitely many weights on the
same prescribed hosts as $u$ grows.  This is strict inclusion of
label sets, not a claim that the corresponding sums have strictly
different numerical values.  The exact local comparison proves
that their difference can be absorbed on the same window.
The remaining all-prime sum, other characters and weight residues,
and $\mathcal R_{\rm cub}$ remain open.  No arbitrary finite
$p$-conductor character or continuous principal-unit character is
declared to occur in $G_u$: its order $144\cdot2^{h(u)}$ is prime
to $269$.  No conclusion that an individual determinant index
vanishes, or that H04.4 is proved, is drawn.

\paragraph{Geometric and source scope.}
The rank-one Artin--Tate coefficients are carried by the existing
finite cover of the arithmetic curve in AC-K25; their stalks and
the full local restrictions are unchanged.  The finite-cover
dictionary is the retained SGA~1, Expos\'e~V, \S\S4 and 7 input.
This proof applies the two original strict theorems only at their
trivial additional finite and continuous characters.  It uses the
complete coefficient/local comparison of AC-K25, the absolute
comparison (FL.4) and codifferent convention of AC-K26-local,
(EU.2), and the full common-window scope of AC-K35/AC-G143.
No new proof of Poitou--Tate duality, an imaginary-quadratic main
conjecture, a full class number, or an implementation theorem is
claimed, and Lean is not executed.
